%% Méthode des rectangles (rectangles « à gauche ») pour approcher l'intégrale de e^{-x}
%% sur [0 ; 1]. À gauche n = 4, à droite n = 10 : la fonction étant décroissante, les
%% rectangles dépassent la courbe et la somme surestime l'intégrale ; l'écart diminue
%% quand n augmente.
\documentclass[tikz,border=4pt]{standalone}
\usepackage{liaisons}
\begin{document}
\begin{tikzpicture}[x=2.9cm,y=2.5cm,
  axe/.style={-{Stealth[length=5pt]}, line width=0.6pt},
  rect/.style={solideB, line width=0.8pt, fill=solideB, fill opacity=0.12},
  courbe/.style={solideA, line width=1.3pt, line join=round},
  grad/.style={font=\scriptsize, inner sep=0pt}]

\def\dx{1.45}   % décalage horizontal du second panneau (en unités d'abscisse)

%% ================= panneau de gauche : n = 4 ==============================
\begin{scope}
  \foreach \i in {0,1,2,3} {
    \pgfmathsetmacro\xa{\i*0.25}
    \pgfmathsetmacro\ha{exp(-\xa)}
    \draw[rect] (\xa,0) rectangle (\xa+0.25,\ha);}
  \draw[axe] (-0.06,0) -- (1.18,0) node[anchor=south east, inner sep=2pt, font=\small] {$x$};
  \draw[axe] (0,-0.05) -- (0,1.18);
  \foreach \xv/\lab in {0.5/{0{,}5}, 1/{1}} {
    \draw[line width=0.5pt] (\xv,0) -- ++(0,-0.06cm);
    \node[grad, anchor=north, yshift=-3pt] at (\xv,0) {$\lab$};}
  \draw[line width=0.5pt] (0,1) -- ++(-0.06cm,0);
  \node[grad, anchor=east, xshift=-3pt] at (0,1) {$1$};
  \node[font=\small, anchor=north east, inner sep=2pt] at (-0.01,-0.01) {$O$};
  \draw[courbe] plot[domain=0:1.1, samples=60, smooth] (\x,{exp(-\x)});
  \node[font=\small, text=solideA, anchor=south west, inner sep=2pt] at (0.72,0.49)
       {$\mathcal{C}_{f}$};
  \node[font=\small, anchor=south] at (0.55,1.2) {$n=4$ : $0{,}714$};
  %% largeur d'un sous-intervalle
  \draw[{Stealth[length=4pt]}-{Stealth[length=4pt]}, black!55, line width=0.6pt]
       (0,-0.2) -- (0.25,-0.2);
  \node[grad, text=black!60, anchor=west, xshift=2pt] at (0.25,-0.2) {$\Delta x=0{,}25$};
\end{scope}

%% ================= filet de séparation ====================================
\draw[black!35, line width=0.4pt] (1.24,-0.2) -- (1.24,1.25);

%% ================= panneau de droite : n = 10 =============================
\begin{scope}[shift={(\dx,0)}]
  \foreach \i in {0,...,9} {
    \pgfmathsetmacro\xa{\i*0.1}
    \pgfmathsetmacro\ha{exp(-\xa)}
    \draw[rect] (\xa,0) rectangle (\xa+0.1,\ha);}
  \draw[axe] (-0.06,0) -- (1.18,0) node[anchor=south east, inner sep=2pt, font=\small] {$x$};
  \draw[axe] (0,-0.05) -- (0,1.18);
  \foreach \xv/\lab in {0.5/{0{,}5}, 1/{1}} {
    \draw[line width=0.5pt] (\xv,0) -- ++(0,-0.06cm);
    \node[grad, anchor=north, yshift=-3pt] at (\xv,0) {$\lab$};}
  \draw[line width=0.5pt] (0,1) -- ++(-0.06cm,0);
  \node[grad, anchor=east, xshift=-3pt] at (0,1) {$1$};
  \node[font=\small, anchor=north east, inner sep=2pt] at (-0.01,-0.01) {$O$};
  \draw[courbe] plot[domain=0:1.1, samples=60, smooth] (\x,{exp(-\x)});
  \node[font=\small, anchor=south] at (0.55,1.2) {$n=10$ : $0{,}664$};
\end{scope}

%% ================= valeur exacte ==========================================
\node[font=\small, anchor=north] at (1.2,-0.33)
     {valeur exacte $\displaystyle\int_{0}^{1}\mathrm{e}^{-x}\,\mathrm{d}x
      =1-\mathrm{e}^{-1}\approx0{,}632$};
\end{tikzpicture}
\end{document}
